% ============================================================================
% COURS 04 — CINÉMATIQUE — PARTIE B
% Source : 5-Cinematique2022.docx
% ============================================================================

\section{Mouvements et trajectoires}
\label{sec:cine-mouvements}

Un mouvement est le déplacement relatif d'un solide observé par rapport
à un solide de référence au cours du temps. La notation $S/0$ désigne
le mouvement du solide $S$ par rapport au solide 0.

\begin{TLCineDef}{Trajectoire}
La trajectoire $\mathcal T_{M\in S/0}$ est le lieu des positions
successives occupées par le point $M$ appartenant au solide $S$, dans
le repère lié au solide 0.
\end{TLCineDef}

Un point de l'espace peut être considéré comme appartenant à un solide
même si aucune matière n'est présente à cet endroit, par exemple un
point situé sur l'axe d'un cylindre creux.

\subsection{Mouvement de translation}

Un solide 2 est en translation par rapport à 1 si :
\[
\boxed{\vec\Omega_{2/1}=\vec0.}
\]
Tous ses points possèdent alors le même vecteur vitesse :
\[
\forall A,B\in2,\qquad
\vec V(A\in2/1)=\vec V(B\in2/1).
\]
On distingue :
\begin{itemize}
  \item la translation rectiligne : trajectoires rectilignes parallèles ;
  \item la translation circulaire : trajectoires circulaires de même rayon ;
  \item la translation curviligne : trajectoires quelconques identiques par translation.
\end{itemize}

\subsection{Mouvement de rotation autour d'un axe fixe}

Le solide 2 est en rotation autour de l'axe $\Delta=(C,\vec u)$ par
rapport à 1 si :
\[
\left\{\mathcal V_{2/1}\right\}_C
=\left\{\begin{array}{c}
\omega\vec u\\\vec0
\end{array}\right\}_C.
\]
Tous les points de l'axe ont une vitesse nulle. Tout point $M$ extérieur
à l'axe décrit un cercle et :
\[
\boxed{\vec V(M\in2/1)=\overrightarrow{MC}\wedge\vec\Omega_{2/1}.}
\]
La norme vaut $V=R|\omega|$, où $R$ est la distance du point à l'axe.

\subsection{Mouvement quelconque}

Lorsqu'un mouvement n'est ni une translation ni une rotation autour d'un
axe fixe, il est qualifié de quelconque. Il peut être analysé comme une
composition de mouvements simples associés aux liaisons successives.

\section{Vecteur rotation et composition}
\label{sec:cine-rotation}

Le vecteur rotation $\vec\Omega_{i/j}$ indique :
\begin{itemize}
  \item la direction de l'axe instantané de rotation ;
  \item la vitesse angulaire en $\mathrm{rad\,s^{-1}}$ ;
  \item le sens positif selon la règle du tire-bouchon.
\end{itemize}

Pour une rotation paramétrée par $\alpha$ autour de $\vec x_1=\vec x_2$ :
\[
\vec\Omega_{2/1}=\dot\alpha\vec x_1.
\]

La composition des rotations s'écrit :
\[
\boxed{\vec\Omega_{n/0}
=\vec\Omega_{n/n-1}+\cdots+\vec\Omega_{2/1}+\vec\Omega_{1/0}.}
\]

\begin{TLCineApp}{Machine cinq axes}
Les solides 1, 4 et 5 sont en translation. Le solide 2 tourne autour de
$\vec x_2$ avec l'angle $\alpha$, et le solide 3 autour de $\vec z_2$
avec l'angle $\beta$.
\[
\vec\Omega_{3/0}
=\dot\beta\vec z_2+\dot\alpha\vec x_2.
\]
Identifier les bases égales et tracer les deux figures de changement de
base.
\end{TLCineApp}

\section{Position, vitesse et accélération}
\label{sec:cine-pva}

\subsection{Vecteur position}

Le vecteur position de $M$ dans le repère
$\mathcal R_0=(O_0,\vec x_0,\vec y_0,\vec z_0)$ est :
\[
\overrightarrow{O_0M}
=x(t)\vec x_0+y(t)\vec y_0+z(t)\vec z_0.
\]
Selon la géométrie, les coordonnées cylindriques ou sphériques peuvent
être plus adaptées.

En coordonnées cylindriques :
\[
\overrightarrow{OM}=r\vec u+z\vec z,
\qquad
\vec u=\cos\theta\vec x+\sin\theta\vec y.
\]
En coordonnées sphériques :
\[
\overrightarrow{OM}=r\vec n,
\qquad
\vec n=\cos\varphi\vec z+\sin\varphi\vec u.
\]

\subsection{Vecteur vitesse}

Le vecteur vitesse est la dérivée du vecteur position dans le repère de
référence :
\[
\boxed{\vec V(M\in S/0)
=\left[\frac{\mathrm d\overrightarrow{O_0M}}{\mathrm dt}\right]_0.}
\]
Il est tangent à la trajectoire et s'exprime en
$\mathrm{m\,s^{-1}}$.

Plus généralement, pour un point fixe $Q$ dans le repère 0 :
\[
\vec V(M\in S/0)
=\left[\frac{\mathrm d\overrightarrow{QM}}{\mathrm dt}\right]_0.
\]

\subsection{Vecteur accélération}

Le vecteur accélération est la dérivée du vecteur vitesse dans le repère
de référence :
\[
\boxed{\vec A(M\in S/0)
=\left[\frac{\mathrm d\vec V(M\in S/0)}{\mathrm dt}\right]_0.}
\]
Il s'exprime en $\mathrm{m\,s^{-2}}$.

\begin{TLCineApp}{Axes de translation d'une machine-outil}
Pour l'extrémité d'outil $P$ :
\[
\overrightarrow{O_0P}
=x(t)\vec x_0+e\vec y_0+(d-z(t)-f)\vec z_0.
\]
On obtient directement :
\[
\vec V(P\in5/0)=\dot x\vec x_0-\dot z\vec z_0,
\]
\[
\vec A(P\in5/0)=\ddot x\vec x_0-\ddot z\vec z_0.
\]
La vitesse maximale vaut
$\sqrt{\dot x_{\max}^2+\dot z_{\max}^2}$.
\end{TLCineApp}

\section{Dérivation vectorielle}
\label{sec:cine-derivation}

Lorsqu'un vecteur est exprimé dans une base mobile, il ne faut pas le
projeter systématiquement dans la base fixe avant de dériver. On utilise
la formule de dérivation vectorielle, dite formule de Boor :

\begin{TLCineDef}{Formule de Boor}
Pour un vecteur $\vec U$ exprimé dans le repère $i$ et dérivé dans le
repère $j$ :
\[
\boxed{
\left[\frac{\mathrm d\vec U}{\mathrm dt}\right]_j
=\left[\frac{\mathrm d\vec U}{\mathrm dt}\right]_i
+\vec\Omega_{i/j}\wedge\vec U.}
\]
Si $\vec U$ est un vecteur unitaire fixe dans $i$ :
\[
\left[\frac{\mathrm d\vec U}{\mathrm dt}\right]_j
=\vec\Omega_{i/j}\wedge\vec U.
\]
\end{TLCineDef}

Pour un repère tournant autour de $\vec x$ avec l'angle $\alpha$ :
\[
\left[\frac{\mathrm d\vec y_2}{\mathrm dt}\right]_0
=\dot\alpha\vec x_2\wedge\vec y_2
=\dot\alpha\vec z_2,
\]
\[
\left[\frac{\mathrm d\vec z_2}{\mathrm dt}\right]_0
=\dot\alpha\vec x_2\wedge\vec z_2
=-\dot\alpha\vec y_2.
\]

\begin{TLCineApp}{Point d'une pièce orientable}
Si
$\overrightarrow{O_0O_3}=y(t)\vec y_0+a\vec z_0+b\vec x_0+c\vec z_2$,
alors :
\[
\vec V(O_3\in3/0)
=\dot y\vec y_0-c\dot\alpha\vec y_2.
\]
En dérivant à nouveau :
\[
\vec A(O_3\in3/0)
=\ddot y\vec y_0-c\ddot\alpha\vec y_2
-c\dot\alpha^2\vec z_2.
\]
Le dernier terme est une accélération centripète.
\end{TLCineApp}
